\documentclass[10pt]{article}

\usepackage[utf8]{inputenc}

\usepackage[T1]{fontenc}

\usepackage{amsmath}

\usepackage{amsfonts}

\usepackage{amssymb}

\usepackage[version=4]{mhchem}

\usepackage{stmaryrd}

\usepackage{mathrsfs}



\begin{document}

уравнений Лапласа по потенциалам мультиполей) следующая формула:



\[

\varepsilon_{\ni Ф}=\varepsilon_{1}\left[1+3 \rho\left(\frac{v+2}{v-1}-\rho+1,3091 \frac{1-v}{v+4 / 3} \rho^{10 / 3}\right)^{-1}\right],

\]



где $v=\varepsilon_{2} / \varepsilon_{1}$, а $\rho$ - объемная плотность тел. В подлинной формуле Рэлея - коэффициент 1,65; правильный коэффициент 1,31 (численное значение суммы нек-рого сходящегося ряда) указан М.К. Поливановым. Если отбросить последнее слагаемое в знаменателе, то получают формулу, выведенную Г. Лоренцом, если два - то формулу, выведенную Дж. Максвеллом. Формулы, полученные Г. Лоренцом и Дж. Максвеллом, учитывают дипольное взаимодействие тел, а формула Рэлея - октупольное.



Учет мультиполей более высокого порядка (см. [5]) дает формула:



\[

\begin{aligned}

& \varepsilon_{\ni \phi}=\varepsilon_{1}\left\{1+3 \rho\left[\frac{v+2}{v-1}-\rho+0,0723 \frac{1-v}{v+6 / 5} \rho^{14 / 3}+0,15256 \frac{1-v}{v+8 / 7} \rho^{6}+\right.\right. \\

&\left.\left.+\frac{1,3091 \frac{1-v}{v+4 / 3} \rho^{10 / 3}\left(1-0,1173 \frac{1-v}{v+6 / 5} \rho^{11 / 3}\right)^{2}}{1+0,4054 \frac{1-v}{v+4 / 3} \rho^{7 / 3}-6,6568 \frac{(1-v)^{2} \rho^{6}}{(v+4 / 3)(v+6 / 5)}}\right]^{-1}\right\}

\end{aligned}

\]



Отбрасывая три слагаемых, пропорциональных $\rho^{11 / 3}$ и $\rho^{6}$ (два), получают формулу, выведенную Р. Мередитом и Ц. Тобиасом [3]. Различие между результатами, полученными по формулам (1) и (2), сушественно для больших плотностей $\rho$ и тем больше, чем больше отличие параметра $v$ от единицы. Формулы (1) и (2) справедливы для тел, размещенных в узлах кубич. решетки, заполняющей все 3-мерное пространство. Если же число «слоев» тел вдоль одной из осей конечно, то $\varepsilon_{\ni Ф}^{(N)}$ возрастает с ростом числа слоев $(N \geqslant 1)$, стремясь к $\varepsilon_{э ф}$ при $N \rightarrow \infty$. Формулы (1) и (2) легко обобшаются на случай кристаллич. решеток других типов (см. [5]).



См. также Области со сложной структурой; краевые задачи.



Лum.: [1] Поливанов М. К., Ферромагнетики, М., 1957; [2] Т олмаче в С.Т., Специальные методы решения задач магнитостатики, Киев, 1983; [3] Meredith R.E., Tobias C.W., «J. Appl. Phys.», 1960, v. 7, р. 1270; [4] Петрин а Д.Я., «Журнал вычислит. математики и математич. физики», 1984, т. 24, с. 709; [5] Малыше в Д. В., Малышев П. В., "Доклады АН УССР», сер. А, 1986, № 2, с. 44; 1987, № 12, с. 48; [6] Математические проблемы теории систем заряженных частиц в неоднородных средах, Киев, 1988, с. 63-76. П. В. Малышев.



ОБЛАСТЬ - см. Топологическое пространство.



ОБЛАСТЬ СУЩЕСТВОВАНИЯ а нал ит иче ской фу н к ц и и - см. Полная аналитическая функция.



ОБЛАСТЬ СХОДИМОСТИ - см. Степенной ряд.



ОБМЕННАЯ ЭНЕРГИЯ - см. Гейзенберга модель.



ОБМЕННОЕ ВЗАИМОДЕЙСТВИЕ - специфическое взаимное влияние одинаковых (тождественных) частиц, эффективно проявляющееся как результат некоторого особого взаимодействия. О. в.- чисто квантовомеханич. Эффект, не имеющий аналога в классич. физике.



Вследствие квантовомеханич. принципа неразличимости одинаковых частиц (тождественности принципа) волновая функция системы должна обладать определенной симметрией относительно перестановки двух таких частиц, то есть их координат и проекций спинов: для частиц с целым спином - бозонов - волновая функция системы не меняется при такой перестановке (является симметричной), а для частиц с полуцелым спином - фермионов - меняет знак (является антисимметричной). Если силы взаимодействия между частищами не зависят от их спинов, волновую функцию системы можно представить в виде произведения двух функций, одна из к-рых зависит только от координат частиц, а другая - только от их спинов. В этом случае из принципа тождественности следует, что координатная часть волновой функции, описывающая движение частиц в пространстве, должна обладать определенной симметрией относительно перестановки координат одинаковых частиц, зависящей от симметрии спиновой части волновой функции. Наличие такой симметрии означает, что имеет место определенная согласованность, корреляция движения одинаковых частиц, к-рая сказывается на энергии системы (даже в отсутствие силовых взаимодействий между частицами). Поскольку обычно влияние частиц друг на друга является результатом действия между ними каких-либо сил, о взаимном влиянии одинаковых частиц, вытекающем из принципа тождественности, говорят как о проявлении специфич. взаимодействия - О. в.



Возникновение О. в. можно проиллюстрировать на примере атома гелия [впервые это было сделано В. Гейзенбергом (W. Heisenberg) в 1926]. Спиновые взаимодействия в легких атомах малы, поэтому волновая функция двух электронов в атоме гелия может быть представлена в виде:



\[

\Psi=\Phi\left(r_{1}, r_{2}\right) \chi\left(s_{1}, s_{2}\right),

\]



где $\Phi\left(\boldsymbol{r}_{1}, \boldsymbol{r}_{2}\right)$ - функция координат электронов, $\chi\left(s_{1}, s_{2}\right)-$ функция проекций их спинов на нек-рое направление. Так как электроны являются фермионами, полная волновая функция должна быть антисимметричной. Если суммарный спин $S$ обоих электронов равен нулю (спины антипараллельны - парагелий), то спиновая функция $\chi$ антисимметрична относительно перестановки спиновых переменных и, следовательно, координатная функция $\Phi$ должна быть симметрична относительно перестановки координат электронов. Если же $S=1$ (спины параллельны - ортогелий), то $\chi$ симметрична, а $\Phi$ антисимметрична. Обозначая через $\psi_{n}\left(r_{1}\right), \psi_{m}\left(r_{2}\right)$ волновые функции отдельных электронов в атоме гелия (индексы $n, m$ означают набор квантовых чисел, определяющих состояние электрона в атоме), можно, пренебрегая сначала взаимодействием между электронами, записать координатную часть волновой функции в виде:



\[

\left.\begin{array}{l}

\Phi_{\mathrm{a}}=\frac{1}{\sqrt{2}}\left[\psi_{n}\left(\boldsymbol{r}_{1}\right) \psi_{m}\left(\boldsymbol{r}_{2}\right)-\psi_{m}\left(\boldsymbol{r}_{1}\right) \psi_{n}\left(\boldsymbol{r}_{2}\right)\right] \text { для } S=1, \\

\Phi_{\mathrm{c}}=\frac{1}{\sqrt{2}}\left[\psi_{n}\left(\boldsymbol{r}_{1}\right) \psi_{m}\left(\boldsymbol{r}_{2}\right)+\psi_{m}\left(\boldsymbol{r}_{1}\right) \psi_{n}\left(\boldsymbol{r}_{2}\right)\right] \text { для } S=0

\end{array}\right\}

\]



(множитель $1 / \sqrt{2}$ введен для нормировки волновой функции). В состоянии с антисимметричной координатной функцией $\Phi_{\mathrm{a}}$ среднее расстояние между электронами оказывается большим, чем в состоянии с симметричной функцией $\Phi_{\mathrm{c}}$; это видно из того, что вероятность $|\Psi|^{2}=\left|\Phi_{\mathrm{a}}\right|^{2}\left|\chi_{\mathrm{c}}\right|^{2}$ нахождения электронов в одной и той же точке $\left(r_{1}=r_{2}\right)$ для состояния $\Phi_{\text {а }}$ равна нулю. Поэтому средняя энергия кулоновского взаимодействия (отталкивания) двух электронов оказывается в состоянии $\Phi_{\mathrm{a}}$ меньшей, чем в состоянии $\Phi_{\mathrm{c}}$. Поправка к энергии системы, связанная с взаимодействием электронов, определяется по теории возмущений:



\[

\mathscr{E}_{\mathrm{B} 3}=K \pm A

\]



где знаки \textbackslash pm относятся соответственно к симметричному и антисимметричному координатным состояниям,



\[

\begin{gathered}

K=e^{2} \int \frac{\left|\psi_{n}\left(\boldsymbol{r}_{1}\right)\right|^{2}\left|\psi_{m}\left(\boldsymbol{r}_{2}\right)\right|^{2}}{\left|\boldsymbol{r}_{1}-\boldsymbol{r}_{2}\right|} d \tau_{1} d \tau_{2}, \\

A=e^{2} \int \frac{\psi_{n}^{*}\left(\boldsymbol{r}_{1}\right) \psi_{m}\left(\boldsymbol{r}_{1}\right) \psi_{m}^{*}\left(\boldsymbol{r}_{2}\right) \psi_{n}\left(\boldsymbol{r}_{2}\right)}{\left|\boldsymbol{r}_{1}-\boldsymbol{r}_{2}\right|} d \tau_{1} d \tau_{2}

\end{gathered}

\]



( $e-$ заряд электрона, $d \tau=d x d y d z-$ элемент объема). Величина $K$ имеет наглядный классич. смысл и соответствует электростатич. взаимодействию двух заряженных «облаков» с плотностями заряда $e\left|\psi_{n}\left(r_{1}\right)\right|^{2}$ и $e\left|\psi_{m}\left(r_{2}\right)\right|^{2}$. Величину $A$, называемую обменным интегралом, можно интерпретировать как электростатич. взаимодействие заряженных «обла-





\end{document}